% Time form: past tense

\newpage
\section{Variant Decision for Thesis Goals}\label{sec:variant-decisions-for-thesis-goals}

During the first sprint of the project, which was focused on project management and planning,
we identified the very promising hypothesis that red and blue team activities could be automated and/or enhanced by \glsxtrshort{AI}.
Pursuing this hypothesis would change the projects focus and goals and thus requires the agreement and approval of the stakeholder.
To communicate this circumstance efficiently and leave the option to go with the status quo, we prepared two options to conduct a \enquote{Variant Decision}.
In the following two subsections the two variants presented in the second checkpoint meeting~\ref{tab:meeting-notes-checkpoint2} are described and illustrated.
It must be noted, that the variants and their visualizations~\ref{fig:variant-1-ai-vs-manually-built-scenarios}\ref{fig:variant-2-multiple-blue-and-red-team-scenarios} are momentary snapshots and will not be updated for traceability reasons.
For example, the Web GUI was already deprioritized during the discussion of the variants, which is reflected in the project goals~\ref{sec:project-goals} but not here.

\subsection{Decision}\label{subsec:variant-decision-for-thesis-goals-decisions}

The stakeholder and project team decided in favor of \textbf{Variant 1: AI vs Manually Built Scenarios (Pivot)}~\ref{subsec:variant-1-ai-vs-manually-built-scenarios} in the checkpoint meeting~\ref{tab:meeting-notes-checkpoint2}.

\subsection{Variant 1: AI vs Manually Built Scenarios (Pivot)}\label{subsec:variant-1-ai-vs-manually-built-scenarios}

Figure~\ref{fig:variant-1-ai-vs-manually-built-scenarios} shows the project goals adjusted for the \glsxtrshort{AI} focused hypothesis, which is favored variant of the project team.
In this option most of the evaluated scenarios will not be realized in favor of \glsxtrshort{AI} based scenarios, which also requires additional research~\ref{subsubsec:ai-integration}.
The basic idea is to implement scenarios for network discovery~\ref{subsubsec:network-discovery-red-team}, intrusion detection~\ref{subsubsec:intrusion-detection-blue-team} and an \glsxtrshort{AI} based scenario with the same capabilities~\ref{subsubsec:ai-purple-team}.
In the evaluation of the thesis we can then assess the quality of the \glsxtrshort{AI} based scenarios itself~\ref{subsubsec:ai-scenario-evaluation} and in comparison to the manually built scenarios~\ref{subsubsec:ai-vs-engineered-scenarios}.

\begin{figure}
    \centering
    \includesvg[width=1.0\linewidth]{../assets/diagrams/planning/planning_variants_1.drawio}
    \caption{Variant 1: AI vs Manually Built Scenarios}
    \label{fig:variant-1-ai-vs-manually-built-scenarios}
\end{figure}

\subsection{Variant 2: Multiple Blue and Red Team Scenarios (Status Quo)}\label{subsec:variant-2-multiple-blue-and-red-team-scenarios}

This is the fallback if the other option~\ref{subsec:variant-1-ai-vs-manually-built-scenarios} is not approved by the stakeholders.
Even if the project team does not favor this more conservative variant, it still represents a valid way to conduct an interesting project.
As illustrated below~\ref{fig:variant-2-multiple-blue-and-red-team-scenarios}, in this variant we propose to proceed with the previously planned scenarios for which we already conducted \enquote{User Story Mappings} documented under the workshops section\ref{sec:workshops}.
The focus would lie in more conventional scenarios, overall framework improvements and an in depth comparison with another framework in form of a GAP Analysis.

\begin{figure}
    \centering
    \includesvg[width=1.0\linewidth]{../assets/diagrams/planning/planning_variants_2.drawio}
    \caption{Variant 2: Multiple Blue and Red Team Scenarios}
    \label{fig:variant-2-multiple-blue-and-red-team-scenarios}
\end{figure}
